\documentclass[11pt]{article}
\usepackage[a4paper,margin=2.2cm]{geometry}
\usepackage[T1]{fontenc}
\usepackage{lmodern}
\usepackage{amsmath,amssymb}
\usepackage{enumitem}
\usepackage{xcolor}
\usepackage[framemethod=default]{mdframed}
\usepackage{hyperref}

\setlist{itemsep=2pt,topsep=3pt}
\newcommand{\done}{\textcolor{green!45!black}{\textsf{[done]}}}
\newcommand{\todo}{\textcolor{orange!85!black}{\textsf{[to do]}}}
\newcommand{\idea}{\textcolor{blue!70!black}{\textsf{[idea]}}}
\newcommand{\open}{\textcolor{red!70!black}{\textsf{[open]}}}
\newcommand{\warn}{\textcolor{red!70!black}{\textsf{[warning]}}}
\newcommand{\Q}[1]{\noindent\textbf{Question.} \emph{#1}\par\smallskip}
\newmdenv[linecolor=gray!60,backgroundcolor=gray!6,innerleftmargin=8pt,innerrightmargin=8pt]{aside}

\title{The story: from FFVL tracks to flying together\\[2pt]\large an outline for the thesis}
\author{Working outline}
\date{\today}

\begin{document}
\maketitle
\setcounter{tocdepth}{1}
\tableofcontents
\bigskip

\noindent Tags: \done{} built and run; \todo{} decided, not built; \idea{} to discuss;
\open{} a question the data have not settled; \warn{} a fact to keep in mind.

\section{Introduction}

\begin{itemize}
  \item Paragliders and hang gliders fly cross-country on thermals; their tracks spread
        superdiffusively. What drives the motion: the route and the relief, the wing, the weather,
        the pilot, flying together?
  \item The thread: clean tracks (2--3), phases (4), groups (5), the factors (6), the route as the
        main driver (7), the group against solo flight (8), a model that reproduces it all (9), the
        tool that lets one see it (10).
  \item Paragliding is the main subject; hang gliders serve for a few comparisons.
\end{itemize}

\section{Dataset}

\begin{itemize}
  \item The FFVL cross-country archive, paragliders (\texttt{para}) and hang gliders
        (\texttt{hang}), IGC files. \done
  \item Since November 2026 the tracks from January 2026 on are 404 on the server: the June copies are
        the only ones. \warn
\end{itemize}

\section{Pre-processing}

\subsection*{The cleaning used today: soaring-anomalous-transport}
\begin{itemize}
  \item Every result today reads the cleaned fixes that soaring-anomalous-transport (tag
        \texttt{legacy-final}) wrote: per segment a uniform grid at the logger interval, smoothed
        $E$, $N$, GNSS altitude, velocities. \done
  \item The procedure is long: duplicates, backward steps, spikes, frozen fixes, altitude bands,
        vertical spikes and level shifts, ground phases trimmed, interior stretches excised, flights
        split into segments, smoothing, and eleven reasons to drop a whole flight.
  \item It drops 15.9\% of the paraglider flights (29\,646 of 186\,052; 20\,798 because no segment
        survived) and 9.3\% of the hang-glider flights (622 of 6\,716).
  \item In the viewer the smoothed track and the raw one differ very little.
\end{itemize}

\subsection*{The cleaning we want: xc-dynamics}
\begin{itemize}
  \item Remove the faulty fixes and flights, and no more; few rules, each easy to explain; fewer
        flights dropped.
  \item Built: altitude channel and barometric witness; top horizontal speed (100 / 175 km/h);
        excursions ($W=20$~s); top vertical speed (30 / 20 m/s); take-off and landing (15 km/h for
        30 s, pauses up to 300 s). \done
  \item Sampling interval: keep the flights logged every 5 s or less, the range where the labelling
        is checked. On the flights kept by soaring this drops 4.3\% of the paragliders (6\,842 of
        156\,406) and 17.8\% of the hang gliders (1\,080 of 6\,094). In the code the limit is still
        10 s. \todo
  \item The procedure is not complete: the remaining steps (gaps and segments next) are still to be
        built. \todo
  \item Then phases, groups and analyses move onto it, and we check that the results hold. \todo
\end{itemize}

\section{Phase labelling}

\Q{At each instant, is the pilot gliding, searching for a thermal or climbing?}
\begin{itemize}
  \item Vilpellet's method: three yes/no features over 30 fixes (straight? rising? turning one
        way?), a three-state hidden Markov model with 8 discrete symbols, Viterbi, a majority vote
        over 90 fixes. Rewritten in xc-lab, same labels bit for bit. \done
  \item Every parameter counts fixes, so the method needs a 1-s grid.
\end{itemize}

\begin{aside}
\textbf{What we do.} Train on the flights logged at 1 Hz, paragliders and hang gliders
separately; apply each model to every flight of its discipline logged every 5 s or less, after
interpolating it onto a 1-s grid with a cubic spline; keep the labels of the recorded fixes.
Today no training of ours exists: every result uses Vilpellet's parameters, fitted on paraglider
tracks at 1 Hz, for both disciplines, and the phases soaring decoded on the 1-Hz flights only.

\smallskip\textbf{Check of the 5-s cadence.} 300 flights logged at 1 Hz, thinned to one fix every
5 s, splined back to 1 s and labelled again: 97.5\% of the fixes keep their 1-Hz label, 96.9\% of
the climbs stay climbs, the climb share is 33.0\% against 33.3\%. \done
\end{aside}

\begin{itemize}
  \item Our own models, trained by Baum--Welch on the 1-Hz flights (109\,364 para, 1\,709 hang), with
        the transition matrix free; compare their labels with Vilpellet's. \todo
  \item Check against the soaring phases used today; port to xc-dynamics. \todo
  \item Quality check against flights segmented by hand (in the viewer, Chapter~10). \todo
  \item Discrete symbols or continuous emissions (turn rate, climb rate)? The latter would be the
        new part of the method. \idea
\end{itemize}

\section{Grouping}

\Q{At every instant of every flight, is the pilot in a group, with whom, and since when?}

\noindent The grouping comes before the factors because the solo population exists only once the
groups are known: Chapters 6 and 7 study solo flights, net of the group.

\begin{itemize}
  \item \textbf{Contact of two pilots.} One is on the path of the other within $\delta=2$~min,
        $\varepsilon_{xy}=500$~m, $\varepsilon_z=300$~m, headings within $\theta_{\max}=45^\circ$ when
        both glide; gaps shorter than $g=120$~s joined; contacts shorter than $\tau^*=38.1$~min
        dropped, the duration above which a time-shift null model explains at most $q=10\%$ of the
        contacts. \done
  \item \textbf{Configuration}: a connected component of the contact graph. \textbf{Group}: the same
        members over a maximal interval. \textbf{Episode}: groups that follow one another and share a
        member. $\tau^*$ applies to each contact, not to membership: two pilots can share a group
        through a third without ever meeting. \done
  \item 21.7\% of the paraglider and 6.9\% of the hang-glider flights fly in a group at some point,
        in 17\,169 episodes. \done
  \item \textbf{Statistics of the groups}: sizes, durations, cohesion, leaders and followers
        (note \texttt{analyses/groups/structure}, to be read and pruned). \done
  \item \textbf{Robustness}: groups and their statistics as $\delta$, $\varepsilon_{xy}$,
        $\varepsilon_z$, $g$, $q$ vary. \todo
  \item Drop the heading condition: the old run already switches it off on flights not at 1 Hz. \idea
  \item The run reads the soaring cleaning; port it to xc-dynamics. \todo
\end{itemize}

\section{Influencing factors}

\Q{How do route and relief, wing, weather and pilot shape the flight statistics, alone and
together?}
\begin{itemize}
  \item On solo flights. Factors: relief (plains, hills, low and high mountains), route, wing class,
        season and hour, weather, pilot.
  \item The factors cannot be assigned, so this is an observational factorial analysis: main effects
        and interactions, with the confounding handled by the design.
  \item Main tool: a mixed model per flight, interactions as product terms, pilot and site-day as
        random effects (they absorb who flies and the weather of the day). Checks: the exact
        decomposition over the $2^k$ cells, which gives main effects and interactions with no model
        (the 8 cells of route, season and hour, wing exist); matched designs. \todo
  \item Overlap: factors that never vary apart (wing and pilot, relief and route) cannot be
        separated. Route as a continuous covariate ($H_{\rm route}$) rather than matched. \open
  \item Weather beyond season and hour needs the soaring-data clouds. \todo
\end{itemize}

\section{Route topology as the main driver of superdiffusion}

\Q{Does the geometry of the route, set by the relief, explain the superdiffusive spreading, and
what does the intermittent component add?}
\begin{itemize}
  \item Geometry first: diffusion against $1/H_{\rm route}$; intermittency on hold. \todo
  \item Plains as the control without ridges. \todo
  \item On solo flights, or solo and group compared? \open
\end{itemize}

\section{Group against solo flight}

\Q{Does flying in a group change how a pilot moves?}

\noindent $S$ solo flights, $G$ group time, $A$ the time alone of the group flights.
\begin{itemize}
  \item \textbf{Raw}: group time is straighter at long lags, climbs more, glides faster. \done
  \item \textbf{Same route} (season, hour, wing): $(G-S|G)-(A-S|A)\approx+0.004$ at
        $H_{\rm loc}(1000~{\rm s})$, because $G-A=+0.037$ is matched by $S|G-S|A=+0.034$: groups fly on
        lines that are straighter for anyone. On a given line the group adds time in climb, not
        straightness. Matching the path nearly fixes the geometry that $H_{\rm loc}$ measures. \done
  \item \textbf{Same flight}: $G-A=+0.037$; entering $+0.027$, leaving $+0.031$. Fixes pilot and day,
        not the line. \done
  \item The two differ by the route: does the group find the line, or the line make the group?
        Test: route-matched controls for the windows before and during the group. \open
  \item Use Chapter 7 to explain the lag beyond which the group effect falls: at long lags the
        ridges bind group and solo pilots alike. Check first which comparison shows the fall. \idea
\end{itemize}

\section{Beyond the PISA model}

\Q{Can one model reproduce the conditional and total MSD, the survivals and the other statistics,
and stay realistic, compatible with what Chapters 4--8 find?}
\begin{itemize}
  \item The PISA model is effective: it reproduces, it does not explain. The aim is a model built
        from the phases, the route and the groups. \todo
  \item On solo flights, or with the group as an ingredient? \open
\end{itemize}

\section{The viewer}

\begin{itemize}
  \item A contribution of its own: tracks, phases and groups on the map, flight by flight. \done
  \item What it shows, how the data reach it, and its uses: checking the cleaning (Chapter 3),
        segmenting flights by hand (Chapter 4), seeing the groups (Chapter 5). \todo
\end{itemize}

\section*{Appendices}
\addcontentsline{toc}{section}{Appendices}
\begin{itemize}
  \item How the uncertainties are assigned: the two-way bootstrap of pilots and blocks. \done
  \item The details of each method that the chapters only cite. \todo
\end{itemize}

\end{document}
